\documentclass[../../main.tex]{subfiles}

\begin{document}

\section{紧性丧失: 无穷维赋范线性空间}\label{section:ch13-紧性丧失}

线性空间 $X$ 称为有限维的, 若存在 $X$ 的子集 $\{e_{1},\cdots,e_{n}\}$ 张成 $X$. 若不存在真子集也张成 $X$, 我们称集合 $\{e_{1},\cdots,e_{n}\}$ 是 $X$ 的一组基, 而称 $n$ 是 $X$ 的维数. 若 $X$ 不能被有限向量族张成, 则称它是无穷维的. 观察到 $X$ 的基 $\{e_{1},\cdots,e_{n}\}$ 在以下意义下是线性无关的:
\begin{align*}
\text{若}\sum_{i=1}^{n}x_{i}e_{i}=0,\quad\text{则对所有 }1\leqslant i\leqslant n,\quad x_{i}=0.
\end{align*}
否则 $\{e_{1},\cdots,e_{n}\}$ 的真子集将张成 $X$.

\begin{theorem}\label{theorem:ch13-有限维范数等价}
有限维线性空间上的任何两个范数等价.
\end{theorem}

\begin{proof}
由于范数的等价是 $X$ 上的范数集合的一个等价关系, 仅须选取 $X$ 上的一个特定范数 $\|\cdot\|_{*}$ 且证明 $X$ 上的任何范数都等价于 $\|\cdot\|_{*}$. 令 $\dim X=n$ 而 $\{e_{1},\cdots,e_{n}\}$ 是 $X$ 的一组基. 对任何 $x=x_{1}e_{1}+\cdots+x_{n}e_{n}\in X$, 定义
\begin{align*}
\|x\|_{*}=\sqrt{x_{1}^{2}+\cdots+x_{n}^{2}}.
\end{align*}
由于欧氏范数是 $\mathbf{R}^{n}$ 上的范数, $\|\cdot\|_{*}$ 是 $X$ 上的范数.

令 $\|\cdot\|$ 为 $X$ 上的任何范数. 我们宣称它等价于 $\|\cdot\|_{*}$. 首先我们找出 $c_{1}\geqslant 0$, 使得
\begin{align}
\text{对所有 }x\in X,\quad\|x\|\leqslant c_{1}\|x\|_{*}.\label{align:ch13-上范数}
\end{align}
事实上, 对 $x=x_{1}e_{1}+\cdots+x_{n}e_{n}\in X$, 根据范数 $\|\cdot\|$ 的次可加性与正齐次性, 以及 $\mathbf{R}^{n}$ 上的 Cauchy-Schwarz 不等式,
\begin{align*}
\|x\|\leqslant\sum_{i=1}^{n}|x_{i}|\|e_{i}\|\leqslant c_{1}\cdot\|x\|_{*},\quad\text{其中 }c_{1}=\sqrt{\sum_{i=1}^{n}\|e_{i}\|^{2}}.
\end{align*}

我们现在找出 $c_{2}>0$, 使得
\begin{align}
\text{对所有 }x\in X,\quad\|x\|_{*}\leqslant c_{2}\cdot\|x\|.\label{align:ch13-下范数}
\end{align}
定义实值函数 $f:\mathbf{R}^{n}\to\mathbf{R}$ 为
\begin{align*}
f(x_{1},\cdots,x_{n})=\left\|\sum_{i=1}^{n}x_{i}e_{i}\right\|.
\end{align*}
该函数是连续的, 由于若把 $\mathbf{R}^{n}$ 考虑为具有欧氏度量的度量空间, 它是具有 Lipschitz 常数 $c_{1}$ 的 Lipschitz 映射. 由于 $\{e_{1},\cdots,e_{n}\}$ 是线性无关的, $f$ 在单位球的边界 $S=\{x\in\mathbf{R}^{n}\mid\sum_{i=1}^{n}x_{i}^{2}=1\}$ 上取正值, 而 $S$ 是紧的因为它是闭与有界的. 紧拓扑空间上的连续实值函数取到最小值. 令 $m>0$ 为 $f$ 在 $S$ 上的最小值. 由范数 $\|\cdot\|$ 的齐次性, 我们得出:
\begin{align*}
\text{对所有 }x\in X,\quad\|x\|\geqslant m\cdot\|x\|_{*}.
\end{align*}
因此\eqref{align:ch13-下范数}对 $c_{2}=1/m$ 成立.
\end{proof}

\begin{corollary}\label{corollary:ch13-有限维同构}
任何两个有相同的有限维数的赋范线性空间是同构的.
\end{corollary}

\begin{proof}
由于同构是赋范线性空间之间的一个等价关系, 仅须证明: 若 $X$ 是维数为 $n$ 的赋范线性空间, 则它同构于欧氏空间 $\mathbf{R}^{n}$. 令 $\{e_{1},\cdots,e_{n}\}$ 为 $X$ 的一组基. 对 $x=(x_{1},\cdots,x_{n})\in\mathbf{R}^{n}$, 定义线性映射 $T:\mathbf{R}^{n}\to X$ 为
\begin{align*}
T(x)=\sum_{i=1}^{n}x_{i}e_{i}.
\end{align*}
由于 $\{e_{1},\cdots,e_{n}\}$ 是基, $T$ 是一对一与映上的. 显然 $T$ 是线性的. 剩下来要证明 $T$ 与它的逆是连续的. 由于线性算子是连续的当且仅当它是有界的, 这相当于证明存在常数 $c_{1}$ 与 $c_{2}$, 使得对每个 $x\in\mathbf{R}^{n}$,
\begin{align*}
\|T(x)\|\leqslant c_{1}\cdot\|x\|_{*}\text{ 且 }\|T(x)\|\geqslant c_{2}\cdot\|x\|_{*}.
\end{align*}
其中 $\|\cdot\|_{*}$ 表示 $\mathbf{R}^{n}$ 的欧氏范数. 这两个常数的存在性从观察到 $x\mapsto\|T(X)\|$ 定义了 $\mathbf{R}^{n}$ 上的范数得出, 这是由于 $\mathbf{R}^{n}$ 上的所有范数是等价的, 都等价于欧氏范数.
\end{proof}

\begin{corollary}\label{corollary:ch13-有限维完备}
任何有限维赋范线性空间是完备的, 因此赋范线性空间的任何有限维子空间都是闭的.
\end{corollary}

\begin{proof}
维数是 $n$ 的有限维空间同构于欧氏空间 $\mathbf{R}^{n}$, 由于 $\mathbf{R}$ 是完备的, 它也是完备的. 由于完备性在同构下保持, 每个有限维赋范线性空间是完备的. 对于赋范线性空间 $X$ 的有限维子空间 $Y$, 由于 $Y$ 具有承袭自范数诱导的度量, 它是完备的, $Y$ 是度量空间 $X$ 的闭子集, 其中 $X$ 为具有由范数诱导的度量的度量空间.
\end{proof}

\begin{corollary}\label{corollary:ch13-有限维闭球紧}
有限维赋范线性空间的闭单位球是紧的.
\end{corollary}

\begin{proof}
令 $X$ 为维数为 $n$ 的赋范线性空间而 $B$ 是它的单位闭球. 令 $T:X\to\mathbf{R}^{n}$ 为同构. 则集合 $T(B)$ 是有界的, 这是由于算子 $T$ 是有界的, 且 $T(B)$ 是闭的由于 $T^{-1}$ 是连续的. 因此, $T(B)$ 作为 $\mathbf{R}^{n}$ 的闭有界子集是紧的. 由于连续映射保持紧性而 $T^{-1}$ 是连续的, $B$ 是紧的.
\end{proof}

\begin{theorem}[Riesz 定理]\label{theorem:ch13-Riesz}
赋范线性空间 $X$ 的闭单位球是紧的当且仅当 $X$ 是有限维的.
\end{theorem}

该定理证明的核心在于以下的引理.

\begin{lemma}[Riesz 引理]\label{lemma:ch13-Riesz引理}
令 $Y$ 为赋范线性空间 $X$ 的闭真线性子空间. 则对每个 $\varepsilon>0$, 存在单位向量 $x_{0}\in X$, 使得
\begin{align*}
\text{对所有 }y\in Y,\quad\|x_{0}-y\|>1-\varepsilon.
\end{align*}
\end{lemma}

\begin{proof}
我们考虑 $\varepsilon=1/2$ 的情形而将一般的情形留作练习. 由于 $Y$ 是 $X$ 的真子集, 我们可以选取 $x\in X\sim Y$. 由于 $Y$ 是 $X$ 的闭子集, 它的补在 $X$ 中是开的, 因此存在与 $Y$ 不交的中心在 $x$ 的球, 即
\begin{align}
\inf\{\|x-y'\|\mid y'\in Y\}=d>0.\label{align:ch13-Riesz距离}
\end{align}
选取向量 $y_{1}\in Y$ 使得
\begin{align}
\|x-y_{1}\|<2d.\label{align:ch13-Riesz逼近}
\end{align}
定义
\begin{align*}
x_{0}=\frac{x-y_{1}}{\|x-y_{1}\|}.
\end{align*}
则 $x_{0}$ 是一个单位向量. 此外, 观察到对任何 $y\in Y$,
\begin{align*}
x_{0}-y=\frac{x-y_{1}}{\|x-y_{1}\|}-y=\frac{1}{\|x-y_{1}\|}\{x-y_{1}-\|x-y_{1}\|y\}=\frac{1}{\|x-y_{1}\|}\{x-y'\},
\end{align*}
其中 $y'=y_{1}+\|x-y_{1}\|y$ 属于 $Y$. 因此, 根据\eqref{align:ch13-Riesz距离}和\eqref{align:ch13-Riesz逼近},
\begin{align*}
\text{对所有 }y\in Y,\quad\|x_{0}-y\|\geqslant\frac{1}{2d}\|x-y'\|\geqslant\frac{1}{2}.
\end{align*}
\end{proof}

\textbf{Riesz 定理的证明.} 我们已证明了有限维赋范线性空间的闭单位球 $B$ 是紧的. 剩下来要证明: 若 $X$ 是无穷维的, 则 $B$ 不是紧的. 假设 $\dim X=\infty$. 我们将归纳地选取 $B$ 中的序列 $\{x_{n}\}$, 使得对 $n\neq m$, $\|x_{n}-x_{m}\|>1/2$. 该序列, 没有 Cauchy 子序列, 因此没有收敛的子序列. 因此 $B$ 不是序列紧的, 由于 $B$ 是一个度量空间, 所以它不是紧的.

剩下来要选取这个序列. 选取任意向量 $x_{1}\in B$. 对于自然数 $n$, 假定我们已选取了 $B$ 中的 $n$ 个向量 $\{x_{1},\cdots,x_{n}\}$. 这些向量的每对距离大于 $1/2$. 令 $X_{n}$ 为这 $n$ 个向量张成的线性空间. 则 $X_{n}$ 是 $X$ 的有限维子空间, 因此是闭的. 此外, $X_{n}$ 是 $X$ 的真子空间由于 $\dim X=\infty$. 根据前一个引理, 我们可以选取 $B$ 中的 $x_{n+1}$, 使得对 $1\leqslant i\leqslant n$, $\|x_{i}-x_{n+1}\|>1/2$. 因此我们已归纳地选取 $B$ 中的序列, 该序列的任何两项距离大于 $1/2$.

\begin{exercise}\label{exercise:ch13-23}
证明有限维赋范线性空间 $X$ 的子集是紧的当且仅当它是闭与有界的.
\end{exercise}

\begin{exercise}\label{exercise:ch13-24}
对 $\varepsilon\neq 1/2$ 完成 Riesz 引理的证明.
\end{exercise}

\begin{exercise}\label{exercise:ch13-25}
给出 $X=\ell^{2}$ 的闭单位球的一个没有有限子覆盖的开覆盖. 接着对 $X=C[0,1]$ 与 $X=L^{2}[0,1]$ 做同样的事情.
\end{exercise}

\begin{exercise}\label{exercise:ch13-26}
对于赋范线性空间 $X$ 和 $Y$, 令 $T:X\to Y$ 是线性的. 若 $X$ 是有限维的, 证明 $T$ 是连续的. 若 $Y$ 是有限维的, 证明 $T$ 是连续的当且仅当 $\ker T$ 是闭的.
\end{exercise}

\begin{exercise}\label{exercise:ch13-27}
(Riesz 定理的另一个证明)令 $X$ 为无穷维赋范线性空间, $B$ 是 $X$ 中的闭单位球, 而 $B_{0}$ 是 $X$ 中的开单位球. 假定 $B$ 是紧的. 则 $B$ 的开覆盖 $\{x+(1/3)B_{0}\}_{x\in B}$ 有有限子覆盖 $\{x_{i}+(1/3)B_{0}\}_{1\leqslant i\leqslant n}$. 用 Riesz 引理(其中 $Y=\operatorname{span}[\{x_{1},\cdots,x_{n}\}]$)导出矛盾.
\end{exercise}

\begin{exercise}\label{exercise:ch13-28}
令 $X$ 为赋范线性空间. 证明 $X$ 是可分的当且仅当存在 $X$ 的紧子集 $K$ 使得 $\overline{\operatorname{span}}[K]=X$.
\end{exercise}

\end{document}